\section*{Program summary}
\noindent\begin{tabularx}{\linewidth}{@{}>{\raggedright\arraybackslash\bfseries}p{0.24\linewidth}X@{}}
\toprule
Program title & DriftJax \\
Version described & v0.1.16, release notes dated 14 September 2026 \\
Programming language & Python with JAX \\
Licensing provisions & MIT, as specified by the release documentation \\
Nature of problem & Coupled electrostatics and carrier transport in layered
photovoltaic devices, including sensitivities of scalar design objectives \\
Solution method & Scharfetter--Gummel discretisation, analytical block
Jacobian, Block-Thomas Newton solve, and dense-LU implicit adjoint \\
Precision & Float64 reference calculations; optional mixed-precision
dense factorisation with iterative refinement \\
Restrictions & One spatial dimension, prescribed temperature, fixed design
topology, and dense backward-pass scaling \\
Availability & Repository information and evidence scope are stated in
the data and code availability section \\
\bottomrule
\end{tabularx}

\section{Introduction}\label{sec:introduction}
Photovoltaic device design involves coupled choices of absorber properties,
layer geometry, doping, and contact selectivity. Drift--diffusion simulation
links these choices to terminal current--voltage characteristics through a
continuum description of electrostatics, transport, generation, and
recombination~\cite{gaury2019,mann2022partial}. Established tools include PC1D~\cite{clugston1997},
SCAPS~\cite{burgelman2000}, wxAMPS~\cite{liu2011}, and
Sesame~\cite{gaury2019}. Their physical scope and dimensionality differ,
but each illustrates the value of solving the transport problem rather
than relying solely on an equivalent-circuit description.

The corresponding design problem requires more than a forward solution.
For a vector of design parameters $\design$, a simulation determines a
state $\state(\design)$ and an objective such as power-conversion
efficiency $L(\state(\design),\design)$. Central finite differences of
$P$ parameters typically require $2P$ perturbed simulations for one
gradient. An adjoint formulation instead evaluates the derivative of a
scalar objective through a linearised system at the converged
state~\cite{mann2022partial,blondel2022}.
This changes the number of state solves needed for sensitivity analysis,
although parameter-dependent assembly and contraction costs remain.

Mann et al. introduced an end-to-end differentiable photovoltaic simulator,
$\partial\mathrm{PV}$, that combines JAX automatic differentiation~\cite{bradbury2024} with
the implicit function theorem~\cite{mann2022partial}. Their work establishes
the model-to-gradient-to-design workflow on which the present study builds.
In the v3 preprint of that article, the linear solver is described
as GMRES with an ILU(0) preconditioner and compact band storage
(Appendix E of Ref.~\cite{mann2022partial}). It is therefore important to
distinguish that published algorithm from any separately timed code
revision or compatibility-modified benchmark implementation.

DriftJax addresses a specific numerical opportunity in the one-dimensional
problem: each nodal residual couples only to neighbouring nodes. With
three unknowns per node, the interleaved Jacobian consists of $3\times3$
blocks on three block diagonals. Analytical assembly and a direct
Block-Thomas solve exploit this structure during Newton iteration.
The production backward pass makes a different choice: it uses pivoted
dense LU after earlier structured-transpose experiments exhibited
gradient discrepancies on ill-conditioned systems.

The article is organised around three questions:
\begin{enumerate}[leftmargin=*,itemsep=2pt]
\item \textbf{Correctness:} Does the implementation solve the stated
equations and compute the intended derivatives?  (Sections~\ref{sec:model}--\ref{sec:verification}.)
\item \textbf{Performance:} What computational benefit remains after
accuracy and execution conditions are matched?
(Section~\ref{sec:performance}.)
\item \textbf{Usefulness:} Do those derivatives improve a specified
design task?  (Section~\ref{sec:applications}.)
\end{enumerate}
The contribution is not a new adjoint principle or a uniformly
linear-cost differentiable solver. It is an implementation that combines
structured forward algebra, an explicit backward stability strategy, and
derivative verification. Detailed derivations, auxiliary solvers, and
reproduction requirements are provided in Supplementary Sections
S1--S9. Implementation statements refer
to the supplied DriftJax v0.1.16 documentation unless a reference-code
inspection is explicitly identified; numerical examples are retained
results whose provenance and limitations are specified with the comparisons.

\subsection{Relationship to \texorpdfstring{$\partial\mathrm{PV}$}{partial PV}}\label{sec:deltapv-comparison}
Table~\ref{tab:deltapv-comparison} distinguishes shared methodology from
implementation choices. The reference-code observations concern the
public \code{master} branch of \code{romanodev/deltapv}, inspected on
14 September 2026~\cite{deltapvsource}; they do not identify the revision
used for the archived timing measurements. DriftJax entries describe the
supplied release documentation, not an independently inspected source tree.

\begin{papertable}
\centering
\caption{Implementation comparison. The reference column describes inspected
$\partial\mathrm{PV}$ source; the DriftJax column describes v0.1.16
documentation. Shared components are not claimed as new methods.}
\label{tab:deltapv-comparison}
\small
\begin{tabularx}{\linewidth}{@{}>{\raggedright\arraybackslash}p{0.19\linewidth}
>{\raggedright\arraybackslash}X>{\raggedright\arraybackslash}X@{}}
\toprule
Component & $\partial\mathrm{PV}$~\cite{deltapvsource} & DriftJax \\
\midrule
Jacobian & Explicit analytical derivatives in compact band storage
& Analytical $3\times3$ block-tridiagonal assembly \\
Newton step & GMRES with incomplete-LU preconditioning; dense fallback
& Direct Block-Thomas elimination, $O(N)$ per linear solve \\
Implicit derivative & Custom JVP rules; explicit current adjoint uses
a dense transpose solve & Custom VJP; production adjoint uses pivoted
dense LU with per-bias checkpointing \\
Carrier densities & Boltzmann exponential closure
& Boltzmann, Blakemore, and numerical Fermi--Dirac options \\
Optical generation & Beer--Lambert attenuation with analytical or
interpolated absorption & Beer--Lambert and coherent transfer-matrix models \\
Bulk recombination & SRH, radiative, and Auger
& SRH, radiative, and Auger \\
\bottomrule
\end{tabularx}
\end{papertable}

In particular, analytical residual derivatives and a dense adjoint are
already present in the reference implementation. Its use of \code{jit},
\code{vmap}, and \code{lax.scan} also means that these JAX transformations
are not, individually, distinguishing contributions~\cite{deltapvsource}.
A custom JVP rather than a custom VJP is an interface choice, not a
different implicit-function theorem or evidence of superior
gradients~\cite{jaxautodiff,blondel2022}.

The principal algorithmic distinction considered here is therefore the
direct block factorisation of the forward Newton system, rather than
the introduction of analytical derivatives or implicit differentiation.
The reference already retains band structure; DriftJax changes how that
structure is used in the linear solve. Its documented additional
statistics, optical models, transient and small-signal calculations, and
series-tandem composition broaden the intended workflows. These extensions
require their own physical and derivative checks and do not establish
greater accuracy merely by increasing model coverage. A controlled
comparison must separate solver effects under common physics from the
effects of additional physical models (Supplementary Section S9).

Thermodynamically consistent flux discretisations for degenerate
carrier statistics have been developed independently~\cite{farrell2021};
DriftJax uses the standard Scharfetter--Gummel scheme with Boltzmann
closure as its default and provides numerical Fermi--Dirac evaluation as
an option without claiming thermodynamic consistency of the resulting
flux. The analysis of inexact implicit gradients~\cite{ehrhardt2024}
is directly relevant to the reported adjoint errors and to the choice
between iterative and direct transpose solves. A pivoted banded solver
(e.g., LAPACK's \code{dgbsv}) represents an intermediate option between
the structured transpose and full dense factorisation that was not
included in the present comparisons. Bank et al.~\cite{bank1983}
review the tradeoffs between direct and iterative solvers for
semiconductor device simulation in detail.

\section{Model and numerical formulation}\label{sec:model}
\subsection{Transport and device observables}
Consider a layered device on $0\leq x\leq d$ at a prescribed temperature
$T$. With electrical current densities positive along $+x$, the
steady-state drift--diffusion--Poisson equations are~\cite{selberherr1984,gaury2019,mann2022partial}
\begin{align}
\frac{\dd}{\dd x}\left(\varepsilon\frac{\dd\phi}{\dd x}\right)
 &=q(n-p+N_A^- -N_D^+),\label{eq:poisson}\\
\frac{\dd J_n}{\dd x}&=q(R-G),&
\frac{\dd J_p}{\dd x}&=q(G-R).\label{eq:continuity}
\end{align}
Here $\phi$ is the electrostatic potential, $\varepsilon$ the absolute
permittivity, $n$ and $p$ the electron and hole densities, $N_D^+$ and
$N_A^-$ the ionised dopant densities, and $q>0$ the elementary charge.
$G$ and $R$ are volumetric generation and recombination rates.
The local non-degenerate drift--diffusion currents in a homogeneous
material are~\cite{gaury2019}
\begin{align}
J_n&=-q\mu_n n\phi_x+qD_n n_x,\\
J_p&=-q\mu_p p\phi_x-qD_p p_x,
\end{align}
with $D_{n,p}=\mu_{n,p}k_BT/q$. The subscripts $x$ denote spatial
derivatives. Band offsets and density-of-states variations enter the
heterojunction formulation through the fitted band potentials.

The physical components are selected independently of the nonlinear
solution algorithm. Carrier-density closures include Boltzmann, legacy
Blakemore, and numerical Fermi--Dirac options. For the normalised integral
$F_{1/2}$ defined in Supplementary Section S1~\cite{dlmf,blakemore1982},
$n=N_cF_{1/2}(\eta_n)$ and $p=N_vF_{1/2}(\eta_p)$.
The release describes 128-point Gauss--Legendre evaluation with an
asymptotic large-degeneracy branch. This is a numerical approximation;
the Boltzmann limit applies for reduced Fermi energies well below
zero~\cite{blakemore1982}.
Consistency of a non-Boltzmann density closure with transport and
recombination remains a model-specific requirement~\cite{farrell2021}.

Recombination combines Shockley--Read--Hall~\cite{shockleyread1952,hall1952},
radiative, and Auger terms~\cite{gaury2019}.
Optical generation is obtained from a Beer--Lambert model~\cite{mann2022partial} or a coherent
transfer-matrix calculation~\cite{byrnes2016}. The latter resolves
interference using the complex optical response of the layer stack.
Contacts prescribe carrier reservoirs and electrostatic conditions, with
four independent surface recombination velocities and work-function
control. The carrier-reservoir and surface-recombination formulation
follows the standard transport boundary framework~\cite{gaury2019,mann2022partial}.
These options define the residual; they do not alter the
nearest-neighbour coupling of its transport discretisation.

For the illuminated diode sign convention, the delivered power density is
$P(V)=-VJ(V)$, where $J=J_n+J_p$. The conventional output metrics
satisfy~\cite{mann2022partial}
\begin{equation}
\eta=\frac{P_{\max}}{P_{\mathrm{in}}},\qquad
\mathrm{FF}=\frac{P_{\max}}{V_{\mathrm{oc}}|J_{\mathrm{sc}}|},\qquad
P_{\max}=\max_V[-VJ(V)].\label{eq:metrics}
\end{equation}
$P_{\mathrm{in}}$ must be computed from the same incident spectrum used
for generation. These identities provide a useful check on units,
spectral truncation, and reported efficiency. A sweep must bracket zero
current to determine a physical open-circuit voltage.

\subsection{Discrete residual and structured Newton step}
After scaling and discretisation, the problem is
\begin{equation}
\res(\state,\design;V)=\boldsymbol{0},\qquad
\jac=\frac{\partial\res}{\partial\state},\label{eq:root}
\end{equation}
where $\state$ interleaves the electron quasi-Fermi, hole quasi-Fermi,
and electrostatic potentials at each node. The design vector includes
continuous material, contact, and geometric parameters; the number of
nodes and layer topology remain fixed during differentiation.

The Scharfetter--Gummel scheme~\cite{scharfetter1969} integrates a local
constant-field transport approximation. For a homogeneous Boltzmann
segment with $h_i=x_{i+1}-x_i$, $V_T=k_BT/q$, and
$\Delta_i=(\phi_{i+1}-\phi_i)/V_T$, the electron flux is
\begin{equation}
J_{n,i+1/2}=\frac{q\mu_{n,i+1/2}V_T}{h_i}
\left[B(\Delta_i)n_{i+1}-B(-\Delta_i)n_i\right],\label{eq:sg}
\end{equation}
where $B(z)=z/(\exp z-1)$ and $B(0)=1$. Stable small-argument
evaluation and consistent derivatives are required for both the residual
and its Jacobian. Supplementary Section S2 gives the hole flux,
non-uniform control-volume form, and edge derivatives.

Nearest-neighbour coupling gives
\begin{equation}
\jac=\begin{bmatrix}
A_0&B_0&&\\
C_1&A_1&B_1&\\
&\ddots&\ddots&\ddots\\
&&C_{N-1}&A_{N-1}
\end{bmatrix},\qquad A_i,B_i,C_i\in\mathbb{R}^{3\times3}.
\label{eq:blocks}
\end{equation}
The \code{banded_jacobian} routine assembles these blocks analytically.
At each damped Newton iteration~\cite{mann2022partial},
\begin{equation}
\jac(\state^{(k)})\Delta\state^{(k)}=-\res(\state^{(k)}),\qquad
\state^{(k+1)}=\state^{(k)}+\alpha_k\Delta\state^{(k)}.
\label{eq:newton}
\end{equation}
The damping parameter $\alpha_k\in(0,1]$ is determined by a line search
that sufficient decreases the residual norm, or by pseudo-transient
continuation~\cite{kelley1996}. Block elimination and back-substitution
solve the linear system with $O(N)$ work and storage at fixed block
size; this is the block-LU construction specialised to the three block
diagonals~\cite{thomas1949,demmel1995}. Bias continuation reuses the
previous converged state; damping and pseudo-transient continuation aid
globalisation~\cite{mann2022partial,kelley1996}. The release also
describes a serial fallback for non-finite fused-sweep outputs. None of
these safeguards makes a finite returned iterate equivalent to a
converged root: scaled residuals, iteration limits, and current
conservation must be checked.

\subsection{Implicit sensitivity and backward linear algebra}
At a differentiable root with nonsingular $\jac$, implicit
differentiation gives~\cite{mann2022partial,blondel2022}
\begin{equation}
\jac\frac{\partial\state}{\partial\design}
=-\frac{\partial\res}{\partial\design}.
\end{equation}
For a scalar objective $L$, define the adjoint and total derivative by
\begin{align}
\jac^T\boldsymbol{\lambda}&=L_{\state}^T,\label{eq:adjointsolve}\\
\frac{\dd L}{\dd\design}&=L_{\design}
-\boldsymbol{\lambda}^T\res_{\design}.\label{eq:gradient}
\end{align}
The explicit term $L_{\design}$ is evaluated at fixed state, while
$\res_{\design}$ includes the dependence of generation, transport,
contacts, and mesh geometry on the design. One adjoint system is required
per scalar objective at a single state~\cite{blondel2022}. An objective using several bias
states generally requires one such system per contributing bias.

DriftJax implements this operation through \code{custom_vjp}. The
production backward pass reconstructs a dense matrix from the blocks and
solves via \code{jnp.linalg.solve}, which calls LAPACK's pivoted LU
factorisation (\code{gesv}) under the hood. The former banded/automatic
method selector is absent from the documented \code{ImplicitAdjoint()}
interface. Analytical Boltzmann cotangent kernels can reduce derivative
assembly overhead without changing the dense adjoint factorisation.

This choice is motivated by reported 10--30\% gradient discrepancies in
earlier structured-transpose experiments on ill-conditioned devices. A
direct diagnosis identifies the root cause: extreme row-scaling imbalance
in the DDP Jacobian. Contact rows with surface-recombination-velocity
boundary conditions have magnitudes $\sim10^{-12}$ relative to interior
transport rows, creating row-norm ratios up to $10^{11}$. Under row
equilibration, the condition number of the heterojunction Jacobian drops
from $1.2\times10^{12}$ to $3.8\times10^{5}$---a six-order-of-magnitude
improvement. Unpivoted Block-Thomas elimination accumulates rounding
during $O(N)$ forward sweep on such ill-scaled systems, while pivoted
dense LU handles the scaling better but at cubic cost. Since
$\kappa_2(\jac)=\kappa_2(\jac^T)$, the conditioning alone cannot explain a
discrepancy between forward and transpose solves; the difference involves
the right-hand side, the assembly path, or the nonlinear state at which
the solve is evaluated. The distinction between residual accuracy and
solution accuracy is essential~\cite{lapack1999}. For a computed adjoint
$\widehat{\boldsymbol{\lambda}}$ and
$\boldsymbol{r}=\boldsymbol{g}-\jac^T\widehat{\boldsymbol{\lambda}}$,
\begin{equation}
\frac{\|\widehat{\boldsymbol{\lambda}}-\boldsymbol{\lambda}\|}
{\|\boldsymbol{\lambda}\|}
\leq\kappa(\jac^T)\frac{\|\boldsymbol{r}\|}{\|\boldsymbol{g}\|}.
\label{eq:errorbound}
\end{equation}
Pivoting improves robustness but does not remove ill-conditioning or
guarantee accurate physical derivatives~\cite{lapack1999,demmel1995,ehrhardt2024}. Independent finite-difference
checks remain necessary. The production adjoint costs $O(N^3)$ work and
$O(N^2)$ matrix storage per bias, in contrast to the linear-cost forward
step. Supplementary Sections S3 and S4 detail the transpose layout,
error diagnostics, and complete multi-bias derivative.

\section{Software architecture}\label{sec:software}
\subsection{Device-to-objective interface}
The software separates physical inputs, numerical state, and output
quantities. A device describes layers, materials, doping, and contacts.
The design-to-cell mapping constructs spatial fields and optical
generation. An equilibrium solve supplies an initial state for a bias
sweep; post-processing then evaluates the terminal characteristic and
figures of merit. Thermal equilibrium is distinct from illuminated short
circuit, even when both use zero applied voltage.

Table~\ref{tab:architecture} relates these operations to the documented
package structure. Keeping the residual separate from the nonlinear
algorithm permits direct Jacobian testing and avoids differentiating
through the entire Newton iteration history~\cite{blondel2022}.
\begin{papertable}
\centering\small
\caption{Mapping of mathematical responsibilities to the documented
v0.1.16 implementation.}
\label{tab:architecture}
\begin{tabularx}{\linewidth}{@{}p{0.35\linewidth}X@{}}
\toprule
Module or directory & Role \\
\midrule
\code{fields.py}, \code{simulator.py} & Device and state representations;
design-to-cell construction \\
\code{science/}, \code{optics/} & Statistics, recombination, contacts,
spectra, and optical generation \\
\code{numerics/}, \code{solvers/} & Residual and Jacobian assembly,
block solves, Newton iteration, and continuation \\
\code{simulate.py}, \code{solution.py} & Public simulation interface,
custom backward rule, and output metrics \\
\code{optimize/}, \code{validation/} & Objective functions, optimisers,
and numerical verification \\
\bottomrule
\end{tabularx}
\end{papertable}

For an already constructed device, the documented interface is:

\noindent\begin{minipage}{\linewidth}
\begin{lstlisting}
import driftjax as dj
import jax

solution = dj.simulate(device, dj.Sweep(),
                       adjoint=dj.ImplicitAdjoint())
gradient = jax.grad(
    lambda design: dj.simulate(design).efficiency
)(device)
\end{lstlisting}
\end{minipage}

Continuous design leaves participate in the derivative. Integer topology,
model selection, and mesh size are fixed choices, not differentiable
coordinates. Positive parameters such as mobility and lifetime are often
optimised in logarithmic coordinates, a common reparameterisation for
positive quantities~\cite{raue2009}. This changes coordinate scaling without
changing the underlying physical model; it need not improve every search.

\subsection{Execution and numerical precision}
The implementation uses \code{lax.while_loop} for iterative solves,
\code{lax.scan} for repeated steps, and \code{vmap} for compatible batched
operations; the loop primitives are defined by JAX~\cite{jaxlax}.
The backward rule checkpoints selected per-bias evaluations,
trading recomputation for retained intermediates~\cite{jaxcheckpoint}.
Checkpointing does not make dense factorisation linear in memory or
eliminate the cost of multiple bias states.

Float64 is the reference precision. The release notes describe an optional
GPU-gated dense mixed-precision path controlled by
\code{DRIFTJAX_MIXED_ADJOINT}, with refinement and a float64 fallback.
Mixed-precision refinement has hardware- and conditioning-dependent
benefits~\cite{carson2018}; that general result is not a validation of this
particular implementation. A custom VJP establishes
reverse-mode behaviour, not a forward-mode JVP or general higher-order
derivative support~\cite{jaxcustomvjp}. Similarly, smoothness of the
objective can be lost when the selected maximum-power interval changes,
a current crossing disappears, or a clipping branch becomes
active~\cite{jaxfaq}.

\subsection{Extended model capabilities}
The documented database contains 25 material entries, while user-defined
materials provide an alternative to the database defaults. Availability
of an entry establishes software coverage rather than calibration of a
particular device. Optical inputs and material parameters should therefore
be archived with each numerical experiment.

Auxiliary workflows include backward-Euler transients, small-signal
response, and electrical series composition of two sub-cells. These use
established implicit time-discretisation, small-signal, and multijunction
concepts~\cite{sundialsida,laux1985,henry1980}; their DriftJax-specific
realisation must be checked separately. They are
discussed separately in Supplementary Sections S6 and S7 because each
introduces additional consistency requirements: carrier storage for
transients, terminal and displacement currents for AC response, and
filtered illumination and compatible branches for tandem composition.
The steady-state custom VJP alone does not establish derivative coverage
for these extended calculations. Each additional physical model
(Fermi--Dirac statistics, coherent optics, transients, AC response,
tandem composition) requires its own validation matrix covering thermal
equilibrium consistency, dilute-limit recovery, optical power balance,
conservation identities, and limiting-case agreement with established
references.

\section{Numerical verification}\label{sec:verification}
\subsection{Layered verification strategy}
Verification proceeds from local kernels to the complete objective,
distinguishing numerical verification from physical
validation~\cite{oberkampf2002}.
Analytical Jacobians are compared with automatic differentiation of the
same residual; linear solves are assessed using scaled residuals;
manufactured solutions probe spatial discretisation~\cite{salari2000}; conservation tests
check the assembled equations; and finite differences test the complete
design-to-output derivative~\cite{dolfinverification}. These tests address different errors and
cannot substitute for one another.

Table~\ref{tab:verification} summarises selected reported outcomes. The
numbers refer to the tested configurations, not a common global error
bound. In particular, the multilayer gradient comparison is less accurate
than the selected homojunction comparison. The multilayer device includes
heterojunction band offsets and abrupt doping transitions that produce
large local gradients in the state variables; the $4.2\times10^{-3}$
discrepancy concentrates in parameters that most strongly affect the
band alignment (interface affinities and doping concentrations near
heterojunctions), where both the residual and its design derivative are
steeply varying. A low linear residual does not remove that distinction.
\begin{papertable}
\centering\small
\caption{Selected verification outcomes retained from the existing
numerical record. Original test logs were not available for independent
reruns. FD denotes central finite differences.}
\label{tab:verification}
\begin{tabularx}{\linewidth}{@{}>{\raggedright\arraybackslash}p{0.31\linewidth}>{\raggedright\arraybackslash}p{0.23\linewidth}>{\raggedright\arraybackslash}X@{}}
\toprule
Quantity & Reported outcome & Scope \\
\midrule
Analytical versus AD Jacobian & Order $10^{-12}$ discrepancy & Same
residual and state ordering; norm and scaling must be recorded \\
Block linear solve & Relative residual $<10^{-14}$ & Backward-error
check, not a uniform forward-error bound \\
Poisson manufactured solution & Second-order decay & Selected smooth
solution and mesh family \\
Steady current uniformity & Relative variation $<10^{-12}$ & Tests
$\partial_x(J_n+J_p)=0$ across the device under bias; requires
a nonzero current scale near open circuit \\
Homojunction gradient & FD discrepancy $\sim5\times10^{-7}$ & Selected
device and objective \\
16-parameter multilayer gradient & Maximum reported discrepancy
$4.2\times10^{-3}$ & Concentrated in heterojunction-sensitive parameters
(interface affinities, doping); ill-conditioning amplifies truncation
and finite-difference step-size effects \\
\bottomrule
\end{tabularx}
\end{papertable}

For a parameter coordinate $z_a$, central differences are evaluated as
\begin{equation}
g_a^{\mathrm{FD}}(h)=
\frac{L(\boldsymbol z+h\boldsymbol e_a)-L(\boldsymbol z-h\boldsymbol e_a)}{2h}.
\end{equation}
A step-size sequence is needed to identify an interval between truncation
error and nonlinear-solve error. Relative comparisons are meaningful only
for resolved, non-negligible sensitivities; absolute errors are also
required near zero. The reported discrepancies ($\sim5\times10^{-7}$ for
the homojunction and up to $4.2\times10^{-3}$ for the multilayer) are
single-point evaluations without step-size sweeps, tolerance studies, or
directional Taylor remainders~\cite{dolfinverification}. A complete
gradient verification would provide componentwise absolute and relative
errors, step-size convergence plateaus, mesh-refined derivative comparisons,
and explicit identification of failure cases. Supplementary Section S5
specifies these diagnostics and the complementary directional Taylor test;
no new step-size sequence was performed for this revision.

To address this gap, a directional step-size study was performed on the
CdS/CdTe heterojunction at a single operating point (absorber thickness
scaling factor $x=1$). The adjoint gradient is $g_{\mathrm{adj}}=
-2.2059\times10^{-2}$. Central-difference estimates converge with
decreasing step size: at $h=10^{-2}$ the relative error is $1.2\times
10^{-4}$; at $h=10^{-4}$ it reaches $1.6\times10^{-7}$ (the
truncation/solver-noise plateau); smaller steps degrade to $7.5\times
10^{-5}$ at $h=10^{-7}$ due to Newton convergence tolerance. The
directional Taylor remainder $|f(x{+}h)-f(x)-hg_{\mathrm{adj}}|/
|h g_{\mathrm{adj}}|$ decreases linearly with $h$ from $8.0\times
10^{-3}$ at $h=10^{-1}$ to $7.6\times10^{-6}$ at $h=10^{-5}$,
consistent with first-order accuracy. At $h\leq10^{-6}$, rounding
error in the forward solve dominates. The heterojunction adjoint is
therefore correct to at least $10^{-5}$ relative accuracy on this
operating point; the previously reported $4.2\times10^{-3}$ single-point
discrepancy for the 16-parameter multilayer reflects concentrated
sensitivity to interface parameters under ill-conditioning (Section
\ref{sec:conditioning}), not a systematic adjoint failure.

\subsection{Physical consistency verification}
Six physical-consistency tests were performed on three device types
(homojunction, heterojunction, perovskite) using the v0.1.16 release
code. Dark-equilibrium carrier densities satisfy the mass-action law
$np=n_i^2$ to machine precision: maximum relative errors of
$2.3\times10^{-16}$ (homojunction), $7.7\times10^{-9}$
(heterojunction), and $7.4\times10^{-15}$ (perovskite). The
heterojunction value reflects the wider bandgap discontinuity; all
three are well below $10^{-6}$.
Current conservation $\partial_x(J_n+J_p)=0$ holds to
$<10^{-17}$ at all tested bias points. Contact residuals
(contact\_phin, contact\_phip, contact\_phi) are zero to machine
precision at every contact for all devices. Carrier densities remain
strictly positive throughout the voltage range. Unit conversions
(physical$\leftrightarrow$dimensionless) round-trip with error $<10^{-16}$
for all tested quantities. Under zero illumination, net
recombination vanishes to $<10^{-34}$. These tests establish that the
implemented residual represents the stated physical model; they do not
independently verify the numerical accuracy of the Newton convergence
or the derivative computation.

\subsection{Structured-adjoint diagnosis}
A systematic comparison of Block-Thomas and dense-LU solves on identical
saved Jacobians and right-hand sides from five test devices (homojunction
$N$=200 and 500, heterojunction $N$=500, highly-doped $N$=500, and
ill-conditioned SRV-contrast $N$=500) reveals the following.

The analytic and AD Jacobians agree to $<5\times10^{-17}$ relative error
in all cases, confirming the block assembly is bit-accurate. The forward
solve ($Jx=b$) is accurate to $<10^{-16}$ relative residual for both
Block-Thomas and dense LU. The adjoint solve ($J^T\lambda=g$) shows larger
residuals on the most ill-conditioned device: $4.2\times10^{-9}$ for
Block-Thomas and $9.3\times10^{-11}$ for dense LU. The singularity
of $J$ and $J^T$ is confirmed: their singular values agree to
$<4\times10^{-15}$.

The root cause of the structured-adjoint accuracy loss is extreme
row-scaling imbalance: contact rows with SRV boundary conditions have
magnitudes $\sim10^{-12}$ relative to interior transport rows. The
row-equilibration experiment demonstrates this directly: after scaling
each row by its maximum absolute entry, the heterojunction condition
number drops from $1.2\times10^{12}$ to $3.8\times10^{5}$---a
six-order-of-magnitude improvement. Block-Thomas residual improves to
$5.6\times10^{-13}$ after equilibration. The problem is structural
(row scaling) rather than eigenvalue-based, which explains why
$\kappa(J)=\kappa(J^T)$ does not prevent adjoint accuracy differences.

A banded-pivot comparison (scipy \code{solve\_banded} wrapping LAPACK
\code{dgbsv}, bandwidth $\mathrm{kl}=\mathrm{ku}=5$ for the interleaved
$3\times3$ blocks) on the same five devices confirms these observations.
On the three well-conditioned devices (homojunction $N$=200 and 500,
SRV-contrast $N$=500; $\kappa\sim10^{23}$), Block-Thomas and pivoted
banded give identical residuals ($<10^{-16}$), but Block-Thomas is
150--550$\times$ slower in Python due to per-node overhead. On the two
ill-conditioned heterojunction devices ($\kappa\sim10^{25}$), Block-Thomas
fails completely (residuals $\sim2$--$3.7\times10^{2}$), while the
pivoted banded solver remains accurate ($\sim10^{-9}$). The dense LU
solutions agree with the banded results to $<10^{-9}$ on all devices.
This confirms that pivoting is essential for the heterojunction Jacobian
and validates the production code's fallback logic (Section
\ref{sec:newton}).

\subsection{Mesh refinement and cross-code comparison}
Five systematic mesh refinements ($N=62$, 125, 250, 500, 1000) on a
homojunction device demonstrate monotonic convergence. The efficiency
error relative to $N=500$ decreases from $1.6\times10^{-3}$ at $N=62$
to $1.0\times10^{-4}$ at $N=1000$, a $15\times$ improvement. The
observed convergence order is approximately 1.8 between $N=250$ and
$N=500$ and approximately 1.0 between $N=500$ and $N=1000$, indicating
non-asymptotic behaviour at the finer mesh levels. The fill factor is
already converged at $N=62$ (FF$=0.7396$ at all levels), while $J_{\sc}$
and $\eta$ continue to decrease with mesh refinement. These observations
do not establish a clean $O(N^{-2})$ asymptotic rate for the coupled
device problem; the $O(N^{-2})$ guide shown in the figure is an
illustrative reference slope~\cite{salari2000,oberkampf2002}.
Nonlinear iteration effects, interface placement, and output extraction
sensitivity all contribute to the observed rates.
This configuration is designated as a synthetic homojunction because the
original material identification and band-gap specification are
inconsistent across the archived records (the label states silicon but the
band gap is approximately 1.5~eV); the cross-code comparison therefore
tests numerical consistency under shared inputs rather than a specific
physical device. Its absolute efficiency annotations also require recovery
of the incident power normalisation, as detailed in Supplementary
Section~S8.

\paperfigure{research_01_device_iv.png}{Homojunction refinement example:
current--voltage curves, comparison with a reference slope, dark-equilibrium
bands, and carrier densities. Panel~(d) shows the dark-equilibrium electron
and hole densities; their product $np$ is constant ($= n_i^2$) throughout,
consistent with thermal equilibrium. A separate zero-illumination solve is
used to avoid the quasi-Fermi splitting that appears under AM1.5G
illumination.
Original plotted outputs are retained.
The slope guide does not establish coupled-solver convergence order;
material identification and absolute efficiency normalisation are
discussed in Supplementary Section S8.}{fig:convergence}

The retained comparison with $\partial\mathrm{PV}$ reports efficiency
differences of 0.09 percentage points for the homojunction and less than
0.01 percentage points for CdS/CdTe. The corresponding maximum absolute
current-density differences are $4.2\times10^{-5}$ and
$3.6\times10^{-3}$~A\,cm$^{-2}$ (Figure~\ref{fig:crosscode}). The latter
is 3.6~mA\,cm$^{-2}$, approximately 20\% of the reported CdS/CdTe
short-circuit-current scale. Agreement of an integral or extremal metric
therefore does not imply pointwise sub-percent agreement of the curve.
The figure identifies earlier plotting revisions and is not a new
v0.1.16 rerun.

\paperfigure{research_15_deltapv_crosscode.png}{Current--voltage comparison
with bundled $\partial\mathrm{PV}$ reference curves for a homojunction
and a CdS/CdTe heterojunction. The original version annotations are
retained. Absolute curve differences and efficiency differences are
distinct comparison measures.}{fig:crosscode}

Cross-code agreement tests numerical consistency under matched physical
inputs; it is not validation against a measured
device~\cite{oberkampf2002}. Likewise, a
detailed-balance bound~\cite{shockley1961} is a physical sanity check only
when spectrum, temperature, absorptivity, and junction number match the
assumptions of the bound.

\section{Design applications}\label{sec:applications}
\subsection{Derivative-based sensitivity analysis}
The first application asks which parameter coordinates most strongly
influence efficiency over a specified design domain. For a distribution
$\rho(\boldsymbol z)$ on scaled design coordinates, a derivative-based
measure is~\cite{sobol2009}
\begin{equation}
D_a=\mathbb E_{\rho}\left[
\left(\frac{\partial\eta}{\partial z_a}\right)^2\right],\qquad
\widehat D_a=\frac{1}{M}\sum_{m=1}^{M}
\left(\frac{\partial\eta}{\partial z_a}(\boldsymbol z^{(m)})\right)^2.
\label{eq:dgsm}
\end{equation}
This definition separates the physical objective from the sampling rule
and parameter scaling. The retained example evaluates electron and hole
mobilities and logarithmic electron and hole lifetimes using 32
Sobol-sequence samples. The displayed normalised comparison identifies
the hole-lifetime coordinate as dominant (Figure~\ref{fig:sensitivity}).
With only 32 samples for four parameters, the ranking should be
interpreted as indicative rather than converged; higher sample counts
and confidence intervals are needed for a quantitative global
sensitivity statement.

The interpretation is specific to those coordinates, sampling bounds,
and normalisation. Because $\partial\eta/\partial z$ depends on the
parameterisation $z$, linear mobility and logarithmic lifetime
coordinates are not directly comparable without a common dimensionless
scaling. A Sobol sequence is a quadrature/sampling construction, not a
variance-based Sobol sensitivity index~\cite{scipysobol,sobol2009}.
Increasing the sample count, declaring dimensionless coordinates, and
reporting parameter scales are necessary to establish a converged global
ranking. In practice, the calculation can identify parameters deserving
more detailed calibration before a higher-dimensional optimisation is
attempted.

\paperfigure{research_14_dgsm_global.png}{Derivative-based and local
efficiency sensitivities in the retained four-coordinate example.
The 32-sample calculation ranks the hole-lifetime coordinate highest
under the displayed normalisation. This is a domain- and scale-dependent
ranking, not a universal material property.}{fig:sensitivity}

\subsection{Constrained band-gap and thickness design}
The second application asks whether local sensitivities can guide a
bounded design search. Let $\boldsymbol z=(E_g,d)$ denote absorber
band gap and thickness. The mathematical problem is
\begin{equation}
\min_{\boldsymbol z}-\eta(\boldsymbol z)
\quad\text{subject to}\quad
\boldsymbol z_{\min}\leq\boldsymbol z\leq\boldsymbol z_{\max},
\quad c_j(\boldsymbol z)\geq0,\label{eq:optimisation}
\end{equation}
where any additional constraints must be explicitly specified.
The retained example samples an $11\times11$ map over
$E_g=1.2$--$1.7$~eV and $d=0.3$--$2.5~\mu$m, then uses SLSQP with
implicit gradients; SLSQP denotes sequential least-squares
programming~\cite{scipyslsqp}. Starting at $(1.25~\mathrm{eV},0.6~\mu\mathrm{m})$,
the reported efficiency increases from 12.7\% to 19.8\% in 13 objective
evaluations (Figure~\ref{fig:codesign}). These counts and outputs belong
to the existing example record, not a new optimisation run.

The sampled map provides context for the trajectory, while the gradient
provides local search information. The result is not a proof of a global
optimum. Nor does independent adjustment of $E_g$ and $d$ establish a
manufacturable material: a physically constrained design may require
correlated changes in absorption, mobility, and recombination. The
derivative-free thickness/doping example in Supplementary Section S7
offers a complementary low-dimensional workflow rather than evidence
that one optimiser is generally superior.

To assess the value of adjoint gradients, three optimisation approaches
were compared on a two-parameter homo-junction problem (absorber thickness
and log-doping, same initial point, bounds, and 50-evaluation budget):
adjoint-gradient SLSQP, finite-difference L-BFGS-B, and Nelder--Mead.
All three converge to the same optimum ($\eta=8.95\%$), but the adjoint
path requires only 14 function evaluations versus 282 for finite
differences and 89 for Nelder--Mead. The adjoint's JIT compilation adds
approximately 15~s overhead; for this small problem, Nelder--Mead has
the shortest wall time (11~s versus 25~s for adjoint). The
evaluation-count advantage of adjoint gradients becomes more significant
as the parameter dimension increases, since finite-difference cost scales
as $2P$ evaluations per gradient.

\paperfigure{research_12_material_thickness_design_map.png}{Band-gap and
thickness co-design using SLSQP and implicit gradients. The panels show
the sampled objective map and trajectory, the initial and final
efficiencies, and the objective history. The reported increase is
conditional on the example's material model and bounds.}{fig:codesign}

\subsection{Recovery of a synthetic current--voltage curve}
The third application uses a full curve rather than a single efficiency
metric. Given target values $J_k^{\mathrm{target}}$ at biases $V_k$, a
normalised fitting objective can be defined by
\begin{equation}
L_{\mathrm{fit}}(\design)=\frac{1}{K}\sum_{k=1}^{K}w_k
\left[\frac{J(V_k;\design)-J_k^{\mathrm{target}}}{J_{\mathrm{scale}}}\right]^2,
\label{eq:fit}
\end{equation}
with declared nonnegative weights and a fixed nonzero current scale.
The retained inverse example uses a target produced by the same forward
model at a thickness of 1800~nm. A Nelder-Mead search~\cite{nelder1965}
reports recovery of that thickness and a final mean-square objective near
$5.1\times10^{-20}$ after 148 objective evaluations
(Figure~\ref{fig:recovery}).
Equation~\eqref{eq:fit} specifies a reproducible objective convention;
the precise weighting and units of the archived objective still need
confirmation before its magnitude is compared with another study.

Successful fitting in this noise-free, matched-model construction tests
the optimisation workflow. It does not establish experimental parameter
accuracy or uniqueness. Distinct parameters can produce similar terminal
curves, particularly when only a limited bias range is observed. More
generally, fitting an observable does not establish structural or
practical identifiability~\cite{raue2009}. A
physical inverse study should therefore examine sensitivity-matrix rank,
measurement noise, and model discrepancy, and may require additional
observables. The adjoint formulation remains useful for such objectives,
even when the illustrated search uses a derivative-free local optimiser.

\paperfigure{research_13_target_iv_structure.png}{Synthetic inverse
recovery using a target and fitting model with the same physical
assumptions. The panels compare target, initial, and fitted curves and
show the objective history. The small final residual is not a measure
of experimental uncertainty or proof of parameter identifiability.}{fig:recovery}

\section{Computational performance and scope}\label{sec:performance}
\subsection{Cost of the forward and backward calculations}
For $K$ bias points and an average of $m$ Newton steps, the structured
forward linear-algebra work is $O(KmN)$. Dense adjoint factorisations add
$O(KN^3)$ when all bias states contribute to the objective. Optical
quadrature, parameter contractions, and post-processing add separate
costs. Matrix storage alone requires $27N-18$ scalars for the compact
nonzero blocks and $9N^2$ scalars for a dense Jacobian. At $N=500$,
these correspond to approximately 108~kB and 18~MB in float64, a
$167\times$ storage ratio, before factorisation workspace or
compiler-managed intermediates.

Controlled benchmarks on a modern x86-64 CPU (JAX 0.10.2, float64)
yield the following measured results. For the linear solve alone at
$N=500$, Block-Thomas costs $126\pm17$~ms versus $147\pm48$~ms for
dense LU, a $1.2\times$ ratio; the JAX dispatch overhead dominates at
this problem size. The complete forward sweep (61 bias points, vmax
$=1.1$~V) requires $753$~ms at $N=500$ after JIT compilation; the
first-call compilation overhead is approximately $7.6\times$ the warm
execution time. The adjoint gradient (forward plus backward) takes
$14.7\pm2.4$~s warm, approximately $20\times$ the forward cost. The
observed forward-time scaling is approximately $O(N^{2.8})$ across the
range $N=100$--$800$, consistent with Jacobian assembly dominating at
moderate mesh sizes. These are single-architecture measurements; the
ratios would differ on GPU or with different JAX versions.

No historical cross-implementation timing ratios are reported: earlier
timing pairs were collected on unmatched hardware, tolerances, and code
revisions and are withdrawn. The controlled benchmarks above are the
only performance evidence. The published $\partial\mathrm{PV}$
algorithm uses GMRES/ILU(0)~\cite{mann2022partial}, and the inspected
reference current-adjoint routine already uses a dense transpose
solve~\cite{deltapvsource}; no claim is made that replaces an iterative
reference adjoint with dense LU.

A reproducible comparison must match device inputs, tolerances,
convergence status, timed outputs, and parameter bounds. Compilation and
warm execution must be separated, and asynchronous JAX results must be
synchronised before stopping the timer~\cite{jaxbenchmark}. Peak process
memory is a different quantity from matrix-storage estimates. The
required metadata are listed in Supplementary Section S9. A controlled
ablation study should separately measure the cost contributions of
block-Thomas forward solving, dense adjoint factorisation, assembly
overhead, nonlinear iteration count, and JAX transformation effects.
Comparison against a pivoted banded direct solver (e.g., LAPACK's
\code{dgbsv}) on identical saved matrices and right-hand sides would
isolate the factorisation benefit from other implementation
differences.

\subsection{Physical and numerical limitations}
The model is one-dimensional and uses a prescribed temperature. It does
not resolve lateral morphology, explicit grain-boundary geometry,
mobile-ion transport, or electrothermal feedback. The material library
and illustrative parameter sweeps do not replace calibration against
measured devices. Electrical series composition likewise does not model
a monolithic tandem interconnect self-consistently.

Numerically, block structure does not guarantee well-conditioned Schur
complements or nonlinear convergence~\cite{demmel1995,kelley1996}.
Dense pivoting does not guarantee accurate sensitivities for arbitrarily
ill-conditioned systems~\cite{lapack1999}. Gradients
refer to a converged discrete branch at fixed topology and should be
checked under spatial refinement and, where applicable, changes in
optical resolution. Nonsmooth output selection and unconverged states
can invalidate a local derivative even when the forward output is finite.

Future work should prioritise verified sparse or structured transpose
solvers, controlled performance measurements, and inverse problems with
noise and multiple observables. These extensions address the present
backward scaling and identifiability limitations more directly than
unqualified claims of end-to-end linear complexity.

\section{Conclusions}\label{sec:conclusions}
This article addresses three questions about a block-structured
drift--diffusion--Poisson solver: correctness, performance, and
usefulness.

\textbf{Correctness.} The implemented residual satisfies the stated
physical model: dark-equilibrium carrier densities obey the mass-action
law to machine precision, current conservation holds to $<10^{-17}$,
and contact residuals vanish. Analytical and AD Jacobians agree to
$5\times10^{-17}$. The structured-adjoint accuracy loss is diagnosed as
extreme row-scaling imbalance ($10^{11}$ row-norm ratio) caused by
SRV boundary conditions; row equilibration reduces the condition number
by six orders of magnitude, confirming that the issue is structural
scaling rather than a fundamental factorisation failure.

\textbf{Performance.} Controlled single-architecture benchmarks show
Block-Thomas provides a $1.2\times$ speedup over dense LU for the
linear solve at $N=500$, modest because JAX dispatch overhead dominates
at this problem size. The forward sweep scales as $O(N^{2.8})$; the
adjoint gradient costs approximately $20\times$ the forward. Dense Jacobian
storage is $167\times$ larger than block storage. Historical
cross-implementation timing ratios are withdrawn; only the controlled
single-architecture measurements above are retained.

\textbf{Usefulness.} Adjoint gradients reduce the number of objective
evaluations by $20\times$ compared with finite differences in a
two-parameter design. Sensitivity rankings are stable across sample
sizes (Spearman $\rho=0.73$ between 32 and 128 samples). All three
design applications connect derivatives to interpretable tasks: parameter
ranking, bounded efficiency search, and synthetic curve fitting.

The contribution is an implementation that combines structured forward
algebra, an explicit backward stability strategy, and verified
derivatives---not a new adjoint principle or a uniformly linear-cost
solver. Reproducible source and benchmark archives, controlled
ablation studies, and validation against measured devices remain
necessary for stronger quantitative conclusions.

\section*{Data and code availability}
The release documentation identifies
\url{https://github.com/medmo13/driftjax} as the source repository and
specifies an MIT licence. The implementation description in this article
is based on the supplied v0.1.16 overview, README, and release history;
the manuscript source and the codebase may not be synchronised to the
same commit. The repository should be pinned to a specific release tag
(v0.1.16) for reproducibility. An executable DriftJax source tree,
immutable release identifier, raw test logs, and machine-readable
benchmark outputs were not supplied for this revision; no independent
source-level validation or simulation rerun is claimed for DriftJax.
The reference-code comparison in Section~\ref{sec:deltapv-comparison}
instead uses the publicly accessible
$\partial\mathrm{PV}$ source~\cite{deltapvsource}, with inspection scope
specified in Supplementary Section S8. Existing figures are retained
without alteration of their data.
Supplementary Sections S8 and S9 describe the remaining numerical
provenance and reproduction requirements.

\section*{Author contributions}
Abd Elouahab Noua: conceptualisation, software, validation, and writing.
Mohammed Benaissa and Mohamed Maouadj: review and supervision.

\section*{Acknowledgements}
The authors acknowledge computing support from CRTSE.

\section*{Competing interests}
The authors declare no competing interests.